\documentclass[12pt,a4paper]{report}
\usepackage[utf8]{inputenc}
\usepackage[T2A]{fontenc}
\usepackage[russian]{babel}
\usepackage{amsmath,amssymb,amsthm,mathtools}
\usepackage{geometry}
\geometry{top=20mm,bottom=20mm,left=25mm,right=15mm}
\usepackage{hyperref}
\hypersetup{colorlinks=true, linkcolor=blue!70!black, citecolor=green!50!black, urlcolor=blue!80!black }
\usepackage{booktabs}
\usepackage{array}
\usepackage{xcolor}
\usepackage{amsopn}
\usepackage{amsfonts}
%\usepackage{unicode-math}
\usepackage{amstext}

% Теоремы, определения, замечания
\theoremstyle{definition}
\newtheorem{definition}{Definition}[chapter]
\newtheorem{example}{Example}[chapter]
\newtheorem{exercise}{Exercise}[chapter]

\theoremstyle{plain}
\newtheorem{theorem}{Theorem}[chapter]
\newtheorem{lemma}[theorem]{Lemma}
\newtheorem{proposition}[theorem]{Proposition}
\newtheorem{corollary}[theorem]{Corollary}

\theoremstyle{remark}
\newtheorem{remark}{Remark}[chapter]

\DeclareMathOperator{\spanop}{span}
\DeclareMathOperator{\rank}{rank}
\DeclareMathOperator{\nullity}{nullity}
\DeclareMathOperator{\im}{im}
\DeclareMathOperator{\diag}{diag}
\DeclareMathOperator{\tr}{tr}
\DeclareMathOperator{\Col}{Col}
\DeclareMathOperator{\Row}{Row}
\DeclareMathOperator{\Null}{Null}

\title{\textbf{\LARGE Линейная алгебра и евклидова геометрия}\\[5mm]
\Large Полное методическое пособие и конспект курса лекций}
\author{\textbf{Курс «Refresher: Линейная алгебра»}}
\date{2026}

\begin{document}
    
    \maketitle
    
    \tableofcontents
    \newpage
    
    \chapter*{Введение}
    \addcontentsline{toc}{chapter}{Введение}
    
    Настоящее методическое пособие предназначено для глубокого и всестороннего освоения базового и продвинутого курсов
    линейной алгебры и евклидовой геометрии, а также для системной подготовки к экзамену.
    
    В основу пособия положен современный \textbf{бескоординатный подход (Map-First Approach)}:
    \begin{enumerate}
        \item \textbf{Первичность геометрических и алгебраических объектов:} Векторы понимаются не как просто
        упорядоченные наборы чисел в таблице, а как элементы абстрактных векторных пространств.
        Линейные отображения рассматриваются первично по отношению к матрицам.
        \item \textbf{Матрицы как координатные представления:} Матрица возникает лишь тогда, когда в пространствах
        прообразов и образов зафиксированы базисы.
        \item Смена базиса изменяет матрицу, но оставляет неизменными фундаментальные инварианты оператора: ранг, след,
        определитель, характеристический полином, собственные и сингулярные значения.
        \item \textbf{Дуальность и измерения:} Строки матрицы интерпретируются как ковекторы (линейные функционалы,
        измерительные приборы), нарезающие пространство на параллельные аффинные гиперплоскости (уровневые множества),
        а столбцы — как образы базисных векторов.
        \item \textbf{Матричные факторизации как декомпозиция структуры:} Разложения $LU$, $QR$, спектральное разложение
        симметричных матриц и сингулярное разложение ($SVD$) представлены как последовательное раскрытие геометрического
        действия оператора в виде суммы простых блоков первого ранга ($u \otimes v^*$).
    \end{enumerate}
    
    Пособие разбито на 4 тематических модуля, полностью соответствующих программе курса и экзаменационным билетам.
    
    \newpage

% =========================================================================
% ГЛАВА 1
% =========================================================================
    
    \chapter{Векторные пространства, дуальность и евклидова геометрия}
    \section{Аксиоматика вещественного векторного пространства}

\begin{definition}[Векторное пространство]
Вещественным \textbf{векторным (линейным) пространством} называется непустое множество $V$, элементы которого называются
\emph{векторами}, на котором заданы две операции:
\begin{enumerate}
    \item \textbf{Сложение векторов:} $+ \colon V \times V \to V$, сопоставляющая каждой паре векторов $u, v \in V$
    вектор $u + v \in V$.
    \item \textbf{Умножение вектора на скаляр:} $\cdot \colon \mathbb{R} \times V \to V$, сопоставляющая скаляру
    $\alpha \in \mathbb{R}$ и вектору $v \in V$ вектор $\alpha v \in V$.
\end{enumerate}
При этом для любых $u, v, w \in V$ и любых $\alpha, \beta \in \mathbb{R}$ выполняются следующие 8 аксиом:
\begin{enumerate}
    \item \textbf{Коммутативность сложения:} $u + v = v + u$;
    \item \textbf{Ассоциативность сложения:} $(u + v) + w = u + (v + w)$;
    \item \textbf{Существование нулевого элемента:} существует элемент $0 \in V$ (называемый \emph{нулевым вектором}),
    что $v + 0 = v$ для любого $v \in V$;
    \item \textbf{Существование противоположного элемента:} для каждого $v \in V$ существует элемент $-v \in V$
    такой, что $v + (-v) = 0$;
    \item \textbf{Унитарность умножения на скаляр:} $1 \cdot v = v$;
    \item \textbf{Ассоциативность умножения на скаляры:} $\alpha (\beta v) = (\alpha \beta) v$;
    \item \textbf{Дистрибутивность относительно сложения векторов:} $\alpha (u + v) = \alpha u + \alpha v$;
    \item \textbf{Дистрибутивность относительно сложения скаляров:} $(\alpha + \beta) v = \alpha v + \beta v$.
\end{enumerate}
\end{definition}

\begin{proposition}[Базовые следствия аксиом]
Из приведённых аксиом строго следуют свойства:
\begin{enumerate}
    \item Нулевой вектор $0 \in V$ единственен.
    \item Для каждого вектора $v$ противоположный вектор $-v$ единственен.
    \item Для любого $v \in V$ выполнено $0 \cdot v = 0_V$.
    \item Для любого $\alpha \in \mathbb{R}$ выполнено $\alpha \cdot 0_V = 0_V$.
    \item $(-1) \cdot v = -v$.
    \item Если $\alpha v = 0$, то либо $\alpha = 0$, либо $v = 0$.
\end{enumerate}
\end{proposition}

\begin{example}[Примеры векторных пространств]
    \begin{enumerate}
        \item Пространство $\mathbb{R}^n$ упорядоченных вещественных $n$-ок чисел со стандартными операциями покомпонентного
        сложения и умножения на скаляр.
        \item Пространство многочленов $\mathbb{P}_n$ степени не выше $n$ от переменной $x$ с вещественными коэффициентами.
        \item Пространство вещественных матриц $\mathbb{R}^{m \times n}$ фиксированного размера с операциями матричного
        сложения и умножения на число.
        \item Пространство непрерывных функций $C[a, b]$ на отрезке.
    \end{enumerate}
\end{example}

\section{Подпространства и линейные комбинации}

\begin{definition}[Subset]
Подмножество $U \subseteq V$ называется \textbf{векторным (линейным) подпространством} пространства $V$ (обозначается $U
\le V$), если $U$ само является векторным пространством относительно операций, индуцированных из $V$.
\end{definition}

\begin{theorem}[Критерий подпространства]
Непустое подмножество $U \subseteq V$ является подпространством тогда и только тогда, когда выполнены следующие условия:
\begin{enumerate}
    \item $0 \in U$ (содержит нулевой вектор);
    \item $\forall u, v \in U \implies u + v \in U$ (\textbf{замкнутость относительно сложения});
    \item $\forall u \in U, \; \forall \alpha \in \mathbb{R} \implies \alpha u \in U$ (\textbf{замкнутость относительно
    умножения на скаляр}).
\end{enumerate}
Эквивалентно: $\forall u, v \in U$ и $\forall \alpha, \beta \in \mathbb{R}$ выполнено $\alpha u + \beta v \in U$.
\end{theorem}

\begin{remark}[Важное отличие: линейные vs аффинные подпространства]
Плоскость или прямая в $\mathbb{R}^3$, не проходящая через начало координат (например, $\{ (x_1, x_2, x_3) \mid x_1 +
x_2 + x_3 = 1 \}$), \textbf{не является} векторным подпространством, поскольку не содержит нулевой вектор $0 = (0,0,0)$
и не замкнута относительно умножения на скаляр ($0 \cdot x = 0 \notin U$).
Такие множества называются \emph{аффинными подпространствами} (сдвинутыми подпространствами).
\end{remark}

\begin{definition}[Линейная комбинация и линейная оболочка]
Пусть $S = \{v_1, v_2, \dots, v_k\} \subset V$.
Вектор вида \[ v = c_1 v_1 + c_2 v_2 + \dots + c_k v_k = \sum_{i=1}^k c_i v_i, \quad c_i \in \mathbb{R}, \]
называется \textbf{линейной комбинацией} векторов множества $S$.

Множество всех всевозможных линейных комбинаций векторов из $S$ называется \textbf{линейной оболочкой} (\textbf{span})
множества $S$ и обозначается:
\[ \spanop(S) = \spanop(v_1, \dots, v_k) = \left\{ \sum_{i=1}^k c_i v_i \;\middle|\; c_i \in \mathbb{R} \right\}. \]
Линейная оболочка любого подмножества $S \subseteq V$ всегда является подпространством в $V$.
Множество $S$ называется \textbf{порождающей системой} (generating set) пространства $U$, если $\spanop(S) = U$.
\end{definition}

\section{Линейная независимость, базис и размерность}

\begin{definition}[Линейная независимость]
Система векторов $\{v_1, v_2, \dots, v_k\} \subset V$ называется \textbf{линейно независимой}, если равенство
\[ c_1 v_1 + c_2 v_2 + \dots + c_k v_k = 0 \]
выполняется тогда и только тогда, когда все коэффициенты тривиальны:
\[ c_1 = c_2 = \dots = c_k = 0. \]
Если же существуют такие коэффициенты $c_1, \dots, c_k$, среди которых хотя бы один отличен от нуля ($ \exists i \colon
c_i \ne 0 $), при которых линейная комбинация равна нулю, система векторов называется \textbf{линейно зависимой}.
\end{definition}

\begin{proposition}[Свойство линейной зависимости]
Система векторов $\{v_1, \dots, v_k\}$ ($k \ge 2$) линейно зависима тогда и только тогда, когда хотя бы один из векторов
системы является линейной комбинацией остальных векторов.
\end{proposition}

\begin{definition}[Базис и размерность]
Упорядоченный набор векторов $\mathcal{B} = (e_1, e_2, \dots, e_n)$ пространства $V$ называется \textbf{базисом}
пространства $V$, если:
\begin{enumerate}
    \item Векторы $e_1, \dots, e_n$ линейно независимы;
    \item Система векторов $e_1, \dots, e_n$ порождает всё пространство $V$, то есть $\spanop(e_1, \dots, e_n) = V$.
\end{enumerate}
Число векторов в базисе называется \textbf{размерностью} пространства $V$ и обозначается $\dim V = n$.
Если базис конечен, пространство называется \emph{конечномерным}.
\end{definition}

\begin{theorem}[О единственности координат в базисе]
Пусть $\mathcal{B} = (e_1, \dots, e_n)$ — базис пространства $V$.
Тогда любой вектор $x \in V$ единственным образом представляется в виде линейной комбинации векторов базиса:
\[ x = x_1 e_1 + x_2 e_2 + \dots + x_n e_n. \]
Числа $x_1, \dots, x_n \in \mathbb{R}$ называются \textbf{координатами} вектора $x$ в базисе $\mathcal{B}$, а
координатный столбец обозначается
$[x]_\mathcal{B} = \begin{pmatrix}
                       x_1 \\ \vdots \\ x_n
\end{pmatrix} \in \mathbb{R}^n$.
\end{theorem}
\begin{proof}
    Существование разложения гарантируется тем, что $\mathcal{B}$ порождает $V$ ($\spanop(\mathcal{B}) = V$).
    Докажем единственность от противного.
    Предположим, что вектор $x$ имеет два различных координатных представления в базисе $\mathcal{B}$:
    \[ x = \sum_{i=1}^n x_i e_i \quad \text{и} \quad x = \sum_{i=1}^n y_i e_i. \]
    Вычтем одно равенство из другого:
    \[ 0 = x - x = \sum_{i=1}^n x_i e_i - \sum_{i=1}^n y_i e_i = \sum_{i=1}^n (x_i - y_i) e_i. \]
    Поскольку по определению базиса векторы $e_1, \dots, e_n$ линейно независимы, линейная комбинация может
    равняться нулю только при равенстве нулю всех ее коэффициентов:
    \[ x_i - y_i = 0 \iff x_i = y_i \quad \forall i \in \{1, 2, \dots, n\}. \]
    Таким образом, наборы коэффициентов тождественно совпадают, что доказывает единственность координат.
\end{proof}

\section{Смена базиса и матрица перехода}

Пусть в пространстве $V$ заданы два базиса: «старый» базис $\mathcal{B} = (e_1, \dots, e_n)$ и «новый» базис
$\mathcal{C} = (e'_1, \dots, e'_n)$.
Каждый вектор нового базиса $e'_j$ можно однозначно разложить по векторам старого базиса:
\[ e'_j = \sum_{i=1}^n p_{ij} e_i = p_{1j} e_1 + p_{2j} e_2 + \dots + p_{nj} e_n. \]
\begin{definition}[Матрица перехода]
Матрица $P = P_{\mathcal{B} \to \mathcal{C}} = [P]_{\mathcal{B} \leftarrow \mathcal{C}} \in \mathbb{R}^{n \times n}$,
столбцами которой являются координаты новых базисных векторов в старом базисе, называется \textbf{матрицей перехода} от
базиса $\mathcal{B}$ к базису $\mathcal{C}$:
\[ P = \begin{pmatrix}
           | & | & & | \\
           [e'_1]_\mathcal{B} & [e'_2]_\mathcal{B} & \dots & [e'_n]_\mathcal{B} \\
           | & | & & |
\end{pmatrix}. \]
Матрица перехода всегда невырождена ($\det P \ne 0$), поскольку ее столбцы линейно независимы.
\end{definition}

\begin{theorem}[Преобразование координат при смене базиса]
Пусть вектор $x \in V$ имеет координаты $[x]_\mathcal{B}$ в базисе $\mathcal{B}$ и координаты $[x]_\mathcal{C}$ в базисе
$\mathcal{C}$.
Тогда:
\[ [x]_\mathcal{B} = P_{\mathcal{B} \leftarrow \mathcal{C}} \, [x]_\mathcal{C}, \quad [x]_\mathcal{C} =
\left(P_{\mathcal{B} \leftarrow \mathcal{C}}\right)^{-1} [x]_\mathcal{B} = P_{\mathcal{C} \leftarrow \mathcal{B}} \,
[x]_\mathcal{B}. \]
\end{theorem}

\section{Дуальные пространства, ковекторы и гиперплоскости}

\begin{definition}[Дуальное пространство и ковектор]
Пусть $V$ — векторное пространство над $\mathbb{R}$.
Линейное отображение $f \colon V \to \mathbb{R}$ называется \textbf{линейным функционалом} (или \textbf{линейной
формой}, или \textbf{ковектором}):
\[ f(\alpha u + \beta v) = \alpha f(u) + \beta f(v) \quad \forall u, v \in V, \; \forall \alpha, \beta \in \mathbb{R}.\]
Множество всех линейных функционалов на $V$ само образует векторное пространство относительно поточечных операций
сложения и умножения на скаляр.
Это пространство называется \textbf{дуальным (сопряжённым)} пространством к $V$ и обозначается $V^* = \mathcal{L}(V,
\mathbb{R})$.
Для конечномерных пространств $\dim V^* = \dim V$.
\end{definition}

\begin{definition}[Дуальный базис]
Пусть $\mathcal{B} = (e_1, \dots, e_n)$ — базис пространства $V$.
Базис $\mathcal{B}^* = (\varepsilon^1, \dots, \varepsilon^n)$ в $V^*$, определяемый условием биортогональности (символ
Кронекера):
\[ \varepsilon^i(e_j) = \delta^i_j = \begin{cases}
                                         1, & i = j, \\ 0, & i \ne j,
\end{cases} \] называется \textbf{дуальным (сопряжённым) базисом}.
\end{definition}

\begin{remark}[Ковекторы как линейные измерительные приборы]
Ковектор $\varepsilon^i$ действует как «измерительный прибор», извлекающий $i$-ю координату вектора.
Если $x = \sum_{j=1}^n x_j e_j$, то:
\[ \varepsilon^i(x) = \varepsilon^i\left(\sum_{j=1}^n x_j e_j\right) = \sum_{j=1}^n x_j \varepsilon^i(e_j) =
\sum_{j=1}^n x_j \delta^i_j = x_i.
\] Следовательно, разложение любого вектора по базису может быть записано через дуальные измерения:
\[ x = \sum_{i=1}^n \varepsilon^i(x) e_i. \]
В матричном представлении вектор $x$ изображается столбцом, а ковектор $f$ — строкой: $f(x) = f x = \sum_i f_i x_i$.
\end{remark}

\begin{proposition}[Геометрия ковектора: нарезка на гиперплоскости]
Ненулевой ковектор $f \in V^*$ нарезает пространство $V$ на семейство параллельных аффинных гиперплоскостей (уровневых
множеств, level sets):
\[ H_c = \{ x \in V \mid f(x) = c \}, \quad c \in \mathbb{R}. \]
При этом:
\begin{enumerate}
    \item При $c = 0$ получаем ядро функционала $\ker f = \{ x \in V \mid f(x) = 0 \}$, которое является векторным
    подпространством коразмерности 1 ($\dim \ker f = n - 1$).
    \item При $c \ne 0$ гиперплоскости $H_c$ параллельны подпространству $\ker f$ и сдвинуты относительно начала
    координат.
\end{enumerate}
\end{proposition}

\section{Иерархия геометрических структур}
    В современной математике евклидова геометрия возникает как результат последовательного обогащения абстрактного
    множества дополнительными математическими структурами:
    \begin{gather*}
        \text{\textbf{Метрическое пространство}} \;\longrightarrow\; \text{\textbf{Нормированное пространство}}
    \;\longrightarrow\;\\
        \text{\textbf{Предгильбертово пространство}} \;\longrightarrow\; \text{\textbf{Евклидово
    пространство}}\\
    \end{gather*}

    \subsection{Метрическое пространство}
    Множество $M$ с функцией расстояния (метрикой) $d \colon M \times M \to \mathbb{R}_{\ge 0}$, удовлетворяющей трём
    аксиомам:
    \begin{enumerate}
        \item Тождество неразличимых: $d(x, y) = 0 \iff x = y$;
        \item Симметрия: $d(x, y) = d(y, x)$;
        \item Неравенство треугольника: $d(x, z) \le d(x, y) + d(y, z)$.
    \end{enumerate}
    В метрическом пространстве нет понятий «направление», «ноль» или «сложение» — есть только расстояния между точками.

    \subsection{Нормированное векторное пространство}
    Векторное пространство $V$, на котором задана \textbf{норма} $\|\cdot\| \colon V \to \mathbb{R}_{\ge 0}$,
    удовлетворяющая аксиомам:
    \begin{enumerate}
        \item Положительная определенность: $\|x\| \ge 0$, причём $\|x\| = 0 \iff x = 0$;
        \item Абсолютная однородность: $\|\alpha x\| = |\alpha| \, \|x\|$ для всех $\alpha \in \mathbb{R}, x \in V$;
        \item Неравенство треугольника (субаддитивность): $\|x + y\| \le \|x\| + \|y\|$.
    \end{enumerate}
    \textbf{Индуцирование метрики:} Всякая норма естественным образом порождает метрику расстояния:
    \[ d(x, y) \coloneqq \|x - y\|. \]
    Обратное неверно: не всякая метрика порождается нормой (например, дискретная метрика $d(x, y) = \begin{cases}
                                                                                                        0, & x=y
                                                                                                        \\ 1, & x \ne y
    \end{cases}$ не удовлетворяет условию однородности).

    \subsection{Пространство со скалярным произведением (предгильбертово)}
    Векторное пространство $V$, оснащенное билинейной операцией $\langle \cdot, \cdot \rangle \colon V \times V \to
    \mathbb{R}$, удовлетворяющей аксиомам:
    \begin{enumerate}
        \item Симметричность: $\langle x, y \rangle = \langle y, x \rangle$;
        \item Линейность по первому аргументу: $\langle \alpha x + \beta y, z \rangle = \alpha \langle x, z \rangle +
        \beta \langle y, z \rangle$;
        \item Положительная определенность: $\langle x, x \rangle \ge 0$, причём $\langle x, x \rangle = 0 \iff x = 0$.
    \end{enumerate}

    \begin{theorem}[Неравенство Коши-Буняковского-Шварца]
Для любых элементов $x, y \in V$ со скалярным произведением справедливо:
\[ |\langle x, y \rangle| \le \sqrt{\langle x, x \rangle} \, \sqrt{\langle y, y \rangle}. \]
    \end{theorem}

    \textbf{Индуцирование нормы:} Всякое скалярное произведение порождает евклидову норму:
    \[ \|x\| \coloneqq \sqrt{\langle x, x \rangle}. \]
    Проверка неравенства треугольника для этой нормы опирается непосредственно на неравенство Коши-Буняковского-Шварца:
    \[ \|x + y\|^2 = \langle x+y, x+y \rangle = \|x\|^2 + 2\langle x, y \rangle + \|y\|^2 \le \|x\|^2 + 2\|x\|\|y\| +
    \|y\|^2 = (\|x\| + \|y\|)^2. \]

    \begin{theorem}[Критерий согласованности нормы со скалярным произведением: тождество параллелограмма]
Норма $\|\cdot\|$ на векторном пространстве $V$ порождается некоторым скалярным произведением тогда и только тогда,
когда для любых $x, y \in V$ выполняется \textbf{тождество параллелограмма}:
\[ \|x + y\|^2 + \|x - y\|^2 = 2\|x\|^2 + 2\|y\|^2. \]
Если тождество параллелограмма выполнено, скалярное произведение однозначно восстанавливается по поляризационному
тождеству:
\[ \langle x, y \rangle = \frac{1}{4}\left(\|x + y\|^2 - \|x - y\|^2\right). \]
    \end{theorem}

    \begin{example}[Контрпример: нормы, не порождаемые скалярным произведением]
Рассмотрим пространство $\mathbb{R}^2$:
\begin{itemize}
    \item Манхэттенская норма ($L_1$): $\|x\|_1 = |x_1| + |x_2|$.
    Возьмём $x = (1, 0)$ и $y = (0, 1)$.
    Тогда $\|x\|_1 = 1$, $\|y\|_1 = 1$.
    Имеем $x+y = (1, 1)$, $x-y = (1, -1)$, отсюда $\|x+y\|_1^2 + \|x-y\|_1^2 = 2^2 + 2^2 = 8$.
    Однако $2\|x\|_1^2 + 2\|y\|_1^2 = 2(1) + 2(1) = 4$.
    Так как $8 \ne 4$, норма $L_1$ не порождается скалярным произведением.
    \item Чебышёвская норма ($L_\infty$): $\|x\|_\infty = \max(|x_1|, |x_2|)$.
    Аналогично не удовлетворяет тождеству параллелограмма.
\end{itemize}
    \end{example}

    \subsection{Евклидово пространство}
    \textbf{Евклидовым пространством} называется конечномерное вещественное векторное пространство, оснащенное скалярным
    произведением.
    В нём одновременно согласованы все уровни геометрии:
    \[ \begin{array}{lcccc}
           \text{Скалярное произведение:} & \langle x, y \rangle & & & \\
           & \Downarrow & & & \\
           \text{Длины (норма):} & \|x\| = \sqrt{\langle x, x \rangle} & & & \\
           & \Downarrow & \searrow & & \\
           \text{Расстояния (метрика):} & d(x, y) = \|x - y\| & & \text{Углы:} & \cos \theta = \frac{\langle x, y
           \rangle}{\|x\| \|y\|}
    \end{array} \]

\section{Ортогональность и ортогональное проецирование}

\begin{definition}[Orthogonality]
Два вектора $u, v \in V$ называются \textbf{ортогональными} ($u \perp v$), если их скалярное произведение равно нулю:
$\langle u, v \rangle = 0$.
\end{definition}

\subsection{Вывод формулы одномерной проекции из ортогональности}
Пусть даны ненулевой вектор $u \in V$ ($u \ne 0$) и произвольный вектор $v \in V$.
Мы ищем ортогональную проекцию вектора $v$ на прямую, натянутую на $u$, то есть вектор $p = P_u(v) \in \spanop(u)$.
\begin{enumerate}
    \item Так как $p \in \spanop(u)$, существует скаляр $\alpha \in \mathbb{R}$ такой, что $p = \alpha u$.
    \item Определим вектор ошибки (перпендикуляр) $e \coloneqq v - p = v - \alpha u$.
    \item Требование \textbf{ортогональности проекции} состоит в том, чтобы вектор ошибки был перпендикулярен
    направляющему вектору $u$:
    \[ e \perp u \iff \langle e, u \rangle = 0 \iff \langle v - \alpha u, u \rangle = 0. \]
    \item Раскрывая скалярное произведение по линейности:
    \[ \langle v, u \rangle - \alpha \langle u, u \rangle = 0 \iff \alpha \langle u, u \rangle = \langle v, u
    \rangle. \]
    \item Поскольку $u \ne 0$, имеем $\langle u, u \rangle = \|u\|^2 > 0$.
    Отсюда однозначно находим коэффициент $\alpha$:
    \[ \alpha = \frac{\langle v, u \rangle}{\langle u, u \rangle} = \frac{\langle v, u \rangle}{\|u\|^2}. \]
    \item Подставляя $\alpha$ в выражение для $p$, получаем каноническую \textbf{формулу ортогональной проекции}:
    \[ P_u(v) = \frac{\langle v, u \rangle}{\langle u, u \rangle} u = \frac{\langle v, u \rangle}{\|u\|^2} u. \]
\end{enumerate}

\subsection{Проекция на многомерное подпространство}
Пусть $U \le V$ — подпространство с базисом, составленным из столбцов матрицы $A\in \mathbb{R}^{n \times k}$ полного
столбцового ранга ($\rank A = k$).
Тогда проекция вектора $y \in \mathbb{R}^n$ на $U = \Col(A)$ представляется в виде $p = A x$ для некоторого $x \in
\mathbb{R}^k$.
Условие ортогональности ошибки $e = y - Ax$ всем столбцам матрицы $A$ записывается в виде:
\[ A^\top (y - Ax) = 0 \iff A^\top A x = A^\top y. \]
Так как столбцы $A$ линейно независимы, матрица Грама $A^\top A \in \mathbb{R}^{k \times k}$ симметрична и строго
положительно определёна, а значит, обратима.
Находим $x$:
\[ x = (A^\top A)^{-1} A^\top y. \]
Тогда проекция равна $p = A x = A (A^\top A)^{-1} A^\top y$.
Отсюда получаем \textbf{матрицу ортогонального проектора}:
\[ P = A (A^\top A)^{-1} A^\top. \]
Матрица $P$ обладает фундаментальными свойствами:
\begin{enumerate}
    \item \textbf{Идемпотентность:} $P^2 = P$ (повторное проецирование ничего не меняет);
    \item \textbf{Симметричность:} $P^\top = P$ (гарантирует ортогональность проецирования).
\end{enumerate}

\section{Метод главных компонент (PCA)}

\textbf{Метод главных компонент (Principal Component Analysis, PCA)} — один из ключевых методов снижения размерности
данных без учителя в машинном обучении и математической статистике.

\subsection{Постановка задачи и алгоритм}
Пусть дана матрица центрированных наблюдений $X \in \mathbb{R}^{n \times d}$, где $n$ — число объектов, $d$ — число
исходных признаков, причём выборочное среднее каждого признака равно нулю ($\sum_{i=1}^n x_{ij} = 0$).
Выборочная ковариационная матрица признаков задаётся формулой:
\[ \Sigma = \frac{1}{n - 1} X^\top X \in \mathbb{R}^{d \times d}. \]
Матрица $\Sigma$ симметрична и положительно полуопределена.

Цель PCA — найти такое ортогональное направление $w_1 \in \mathbb{R}^d$ ($\|w_1\|_2 =1$), проекция данных на которое
обладает максимальной выборочной дисперсией:
\[ \max_{\|w_1\| = 1} \operatorname{Var}(X w_1) = \max_{\|w_1\| = 1} \frac{1}{n-1} (X w_1)^\top (X w_1) =
\max_{\|w_1\| = 1} w_1^\top \Sigma w_1. \]
Используя метод множителей Лагранжа для функции $L(w_1, \lambda) =w_1^\top \Sigma w_1 - \lambda (w_1^\top w_1 - 1)$,
получаем необходимое условие экстремума:
\[ \nabla_{w_1} L = 2 \Sigma w_1 - 2 \lambda w_1 = 0 \iff \Sigma w_1 = \lambda w_1. \]
Таким образом, направление первой главной компоненты $w_1$ — это \textbf{собственный вектор} ковариационной матрицы
$\Sigma$, отвечающий ее максимальному собственному значению $\lambda_1$.
Дисперсия вдоль этого направления равна в точности $\lambda_1$.

Последующие главные компоненты $w_2,\dots,w_k$ находятся как собственные векторы $\Sigma$, ортогональные предыдущим,
упорядоченные по убыванию собственных значений $\lambda_1 \ge \lambda_2 \ge \dots \ge \lambda_d \ge 0$.

\subsection{Связь PCA с сингулярным разложением (SVD)}
На практике вычисление ковариационной матрицы $\Sigma = \frac{1}{n-1} X^\top X$ не требуется.
Достаточно применить SVD непосредственно к центрированной матрице данных:
\[ X = U S V^\top. \]
Тогда:
\[ X^\top X = (V S U^\top)(U S V^\top) = V S^2 V^\top. \]
Отсюда видно, что:
\begin{enumerate}
    \item Правые сингулярные векторы $V = (v_1, \dots, v_d)$ являются в точности главными компонентами (собственными
    векторами матрицы ковариации $\Sigma$).
    \item Собственные значения ковариационной матрицы связаны с сингулярными числами матрицы $X$ соотношением:
    \[ \lambda_i = \frac{s_i^2}{n - 1}. \]
    \item Новые низкоразмерные координаты данных (проекции) вычисляются умножением: $Z_k = X V_k = U_k S_k$.
\end{enumerate}

\newpage

% =========================================================================
% ГЛАВА 2
% =========================================================================

    \chapter{Матрицы, линейные отображения, ранг и системы уравнений}
    \section{Линейные отображения и их координатное представление}

\begin{definition}[Линейное отображение]
Отображение $T \colon V \to W$ между вещественными векторными пространствами называется \textbf{линейным отображением}
(линейным оператором / гомоморфизмом), если для любых $u, v \in V$ и любого $\alpha \in \mathbb{R}$ выполнены:
\begin{enumerate}
    \item Аддитивность: $T(u + v) = T(u) + T(v)$;
    \item Однородность: $T(\alpha u) = \alpha T(u)$.
\end{enumerate}
Эквивалентно сохранению произвольных линейных комбинаций:
\[ T\left(\sum_{i=1}^k c_i v_i\right) = \sum_{i=1}^k c_i T(v_i). \]
\end{definition}

\begin{definition}[Ядро и образ]
Пусть $T \colon V \to W$ — линейное отображение.
\begin{itemize}
    \item \textbf{Ядром} (\textbf{kernel / nullspace}) отображения $T$ называется множество векторов из $V$,
    переходящих в нулевой вектор пространства $W$: \[ \ker(T) = \{ v \in V \mid T(v) = 0_W \} \le V \].
    \item \textbf{Образом} (\textbf{image / range}) отображения $T$ называется подмножество векторов из $W$, которые
    достигаются отображением $T$:
    \[ \im(T) = \{ w \in W \mid \exists v \in V \colon T(v) = w \} \le W. \]
\end{itemize}
Отображение $T$ \textbf{инъективно} тогда и только тогда, когда $\ker(T) = \{0_V\}$.\\ Отображение $T$
\textbf{сюръективно} тогда и только тогда, когда $\im(T) = W$.
\end{definition}

\section{Матрица линейного отображения: почему столбцы суть образы базисных векторов?}

Пусть в пространстве $V$ зафиксирован базис $\mathcal{B}=(v_1,\dots, v_n)$, а в пространстве $W$ зафиксирован базис
$\mathcal{C} = (w_1, \dots, w_m)$.

Возьмём произвольный вектор $x \in V$.
Разложим его по базису $\mathcal{B}$:
\[ x = x_1 v_1 + x_2 v_2 + \dots + x_n v_n = \sum_{j=1}^n x_j v_j. \]
Применим линейное отображение $T$ к вектору $x$, используя свойство линейности:
\[ T(x) = T\left(\sum_{j=1}^n x_j v_j\right) = \sum_{j=1}^n x_j T(v_j). \]
Отсюда виден ключевой фундаментальный факт: \textbf{действие линейного отображения на любой вектор полностью и
однозначно определяется тем, куда переходят базисные векторы} $T(v_1), \dots, T(v_n) \in W$.

Теперь разложим каждый вектор $T(v_j) \in W$ по базису $\mathcal{C}$:
\[ T(v_j) = \sum_{i=1}^m a_{ij} w_i = a_{1j} w_1 + a_{2j} w_2 + \dots + a_{mj} w_m. \]
Координатный столбец вектора $T(v_j)$ в базисе $\mathcal{C}$ равен:
\[ [T(v_j)]_\mathcal{C} = \begin{pmatrix}
                              a_{1j} \\ a_{2j} \\ \vdots \\ a_{mj}
\end{pmatrix}. \]
Подставим это разложение в выражение для $T(x)$:
\[ T(x)=\sum_{j=1}^n x_j\left(\sum_{i=1}^m a_{ij} w_i\right)=\sum_{i=1}^m \left( \sum_{j=1}^n a_{ij} x_j\right) w_i.
\]
Следовательно , $i$-я координата вектора $y = T(x)$ в базисе $\mathcal{C}$ равна $y_i = \sum_{j=1}^n a_{ij} x_j$.
В матричной форме это записывается как:
\[ \begin{pmatrix}
       y_1 \\ \vdots \\ y_m
\end{pmatrix} = \begin{pmatrix}
                    a_{11} & a_{12} & \dots & a_{1n} \\ a_{21} & a_{22} & \dots & a_{2n} \\ \vdots
                    & \vdots & \ddots & \vdots \\ a_{m1} & a_{m2} & \dots & a_{mn}
\end{pmatrix} \begin{pmatrix}
                  x_1 \\ \vdots \\ x_n
\end{pmatrix} \iff [T(x)]_\mathcal{C} = [T]_{\mathcal{C} \leftarrow
\mathcal{B}} [x]_\mathcal{B}. \]

\begin{theorem}[Структура матрицы линейного оператора]
Матрица линейного оператора $[T]_{\mathcal{C} \leftarrow \mathcal{B}} \in \mathbb{R}^{m \times n}$ состоит из столбцов,
каждый $j$-й столбец которых является координатным представлением образа $j$-го базисного вектора $v_j$ в базисе
$\mathcal{C}$:
\[ [T]_{\mathcal{C} \leftarrow \mathcal{B}} = \begin{pmatrix}
                                                  | & | & & | \\ [T(v_1)]_\mathcal{C} & [T(v_2)]_\mathcal{C}
                                                  & \dots & [T(v_n)]_\mathcal{C} \\
                                                  | & | & & |
\end{pmatrix}. \]
\end{theorem}

\section{Концептуальный смысл теоремы о ранге и дефекте (Rank-Nullity)}

\begin{theorem}[Rank-Nullity Theorem]
Пусть $V$ и $W$ — конечномерные векторные пространства, и $T \colon V \to W$ — линейное отображение.
Тогда:
\[ \dim V = \dim(\ker T) + \dim(\im T). \]
В терминах ранга и дефекта (дефект $\nullity(T) \coloneqq \dim(\ker T)$, ранг $\rank(T) \coloneqq \dim(\im T)$):
\[ \dim V = \rank(T) + \nullity(T). \]
Для матрицы $A \in \mathbb{R}^{m \times n}$ с $n$ столбцами:
\[ n = \rank(A) + \dim(\Null A). \]
\end{theorem}

\begin{proof}[Идея доказательства]
Пусть $\dim V = n$, и пусть $\dim(\ker T) = k$.
Выберем базис ядра $\ker T$: пусть это векторы $(u_1, u_2, \dots, u_k)$.
Поскольку эти векторы линейно независимы в $V$, по теореме о дополнении базиса их можно дополнить векторами $(v_1,
\dots, v_{n-k})$ до полного базиса всего пространства $V$:
\[ \mathcal{B} = (u_1, \dots, u_k, v_1, \dots, v_{n-k}). \]
Рассмотрим образы дополнительных векторов в $W$: $w_j \coloneqq T(v_j)$ для $j = 1, \dots, n-k$.
Покажем, что система $(w_1, \dots, w_{n-k})$ образует базис подпространства $\im T$:
\begin{enumerate}
    \item \textbf{Порождаемость:} Любой вектор $x \in V$ раскладывается как $x = \sum_{i=1}^k \alpha_i u_i +
    \sum_{j=1}^{n-k} \beta_j v_j$.
    Применяя $T$ и учитывая, что $T(u_i) = 0$:
    \[ T(x) = \sum_{i=1}^k \alpha_i T(u_i) + \sum_{j=1}^{n-k} \beta_j T(v_j) = 0 + \sum_{j=1}^{n-k} \beta_j w_j. \]
    Следовательно, векторы $w_1, \dots, w_{n-k}$ порождают $\im T$.
    \item \textbf{Линейная независимость:} Пусть $\sum_{j=1}^{n-k} c_j w_j = 0$.
    Тогда $T\left(\sum_{j=1}^{n-k} c_j
    v_j\right) = 0$, что означает, что вектор $v \coloneqq \sum_{j=1}^{n-k} c_j v_j$ лежит в ядре $\ker T$.
    Но тогда он должен выражаться через базис ядра $(u_1, \dots, u_k)$:
    \[ \sum_{j=1}^{n-k} c_j v_j = \sum_{i=1}^k d_i u_i \iff \sum_{j=1}^{n-k} c_j v_j - \sum_{i=1}^k d_i u_i = 0. \]
    Так как векторы $(u_1,\dots, u_k,v_1,\dots, v_{n-k})$ образуют базис $V$, они линейно независимы, следовательно, все
    коэффициенты $c_j = 0$ и $d_i = 0$.
\end{enumerate}
Итак, $(w_1, \dots, w_{n-k})$ — базис $\im T$, а значит, $\dim(\im T) = n - k$.
Складывая: $\dim(\ker T) + \dim(\im T) = k + (n - k) = n = \dim V$.
\end{proof}

\begin{remark}[Концептуальная геометрическая интерпретация]
Теорема о ранге и дефекте выражает \textbf{закон сохранения размерности}: Исходное пространство размерности $n$ при
линейном отображении делится на две непересекающиеся составляющие:
\begin{itemize}
    \item Направления, которые отображение \textbf{схлопывает в ноль} (ядро $\ker T$ размерности $\nullity T$).
    \item Направления, которые \textbf{переживают трансформацию} и разворачиваются в образе
    ($\im T$ размерности $\rank T$).
\end{itemize}
Ни одно измерение не может исчезнуть бесследно: либо степень свободы уничтожается отображением (уходит в дефект), либо
проявляется на выходе (даёт вклад в ранг).
\end{remark}

\section{Четыре фундаментальных подпространства матрицы}

Для матрицы $A \in \mathbb{R}^{m \times n}$ канонически определены 4 подпространства (теорема Странга):
\begin{enumerate}
    \item \textbf{Пространство столбцов} (\textbf{Column space}): $\Col(A) = \im(A) \le \mathbb{R}^m$,
    размерность $\rank(A) = r$;
    \item \textbf{Пространство строк} (\textbf{Row space}): $\Row(A) = \Col(A^\top) \le \mathbb{R}^n$, размерность
    $\rank(A^\top) = \rank(A) = r$;
    \item \textbf{Нуль-пространство (ядро)} (\textbf{Nullspace}): $\Null(A) = \ker(A) \le \mathbb{R}^n$,
    размерность $n - r$;
    \item \textbf{Левое нуль-пространство}: $\Null(A^\top) = \ker(A^\top) \le \mathbb{R}^m$, размерность $m - r$.
\end{enumerate}

\begin{theorem}[Ортогональные дополнения фундаментальных подпространств]
В пространствах $\mathbb{R}^n$ и $\mathbb{R}^m$ фундаментальные подпространства взаимно ортогональны и образуют полные
прямые суммы:
\begin{gather*}
    \Null(A) = (\Row(A))^\perp \quad \text{в} \quad \mathbb{R}^n \implies \mathbb{R}^n = \Row(A) \oplus \Null(A),\\
    \Null(A^\top) = (\Col(A))^\perp \quad \text{в} \quad \mathbb{R}^m \implies \mathbb{R}^m = \Col(A) \oplus
    \Null(A^\top).\\
\end{gather*}
\end{theorem}

\section{Диагностика невырожденности: определитель, ранг и обратимость}

Для квадратной матрицы $A\in\mathbb{R}^{n\times n}$ понятия определителя, ранга и обратимости представляют собой три
взаимодополняющих взгляда на одно и то же явление — \textbf{невырожденность оператора}:

\begin{theorem}[Теорема об эквивалентности условий невырожденности]
Для квадратной матрицы $A \in \mathbb{R}^{n \times n}$ следующие утверждения строго эквивалентны:
\begin{enumerate}
    \item Матрица $A$ обратима (существует $A^{-1}$ такая, что $A A^{-1} = A^{-1} A = I_n$);
    \item Определитель отличен от нуля: $\det(A) \ne 0$;
    \item Матрица $A$ имеет полный ранг: $\rank(A) = n$;
    \item Ядро матрицы тривиально: $\Null(A) = \{0\}$;
    \item Столбцы матрицы $A$ линейно независимы и образуют базис $\mathbb{R}^n$;
    \item Строки матрицы $A$ линейно независимы и образуют базис $\mathbb{R}^n$;
    \item Система линейных уравнений $Ax = 0$ имеет исключительно тривиальное решение $x = 0$;
    \item Для любого вектора $b \in \mathbb{R}^n$ система $Ax = b$ имеет единственное решение $x = A^{-1} b$;
    \item Число 0 не является собственным значением матрицы $A$ ($\lambda_i \ne 0$ для всех $i$);
    \item Все сингулярные числа матрицы $A$ строго положительны ($\sigma_n > 0$).
\end{enumerate}
\end{theorem}

\textbf{Геометрическая интерпретация проблемы:}
\begin{itemize}
    \item $\det(A) = 0$ означает, что оператор сжимает $n$-мерный объём единичного гиперкуба в нуль
    (отображение сплющивает пространство в подпространство меньшей размерности).
    \item $\rank(A) < n$ указывает размерность полученного сплющенного образа.
    \item Несуществование $A^{-1}$ констатирует невозможность обратить этот процесс вспять
    (информация о направлении, попавшем в ядро, безвозвратно утеряна).
\end{itemize}

\section{Системы линейных уравнений и элементарные операции}

Система $m$ линейных уравнений с $n$ неизвестными записывается в виде:
\[ Ax = b, \quad A \in \mathbb{R}^{m \times n}, \; x \in \mathbb{R}^n, \; b \in \mathbb{R}^m. \]
\begin{theorem}[Теорема Кронекера-Капелли]
Система линейных уравнений $Ax = b$ совместна (имеет хотя бы одно решение) тогда и только тогда, когда ранг основной
матрицы равен рангу расширенной матрицы:
\[ \rank(A) = \rank(A \mid b). \]
При этом:
\begin{enumerate}
    \item Если $\rank(A) = \rank(A \mid b) = n$ (число неизвестных), система имеет \textbf{ровно одно решение}.
    \item Если $\rank(A) = \rank(A \mid b)=r < n$, система имеет \textbf{бесконечно много решений}, зависящих от $n - r$
    свободных параметров.
    \item Если $\rank(A) < \rank(A \mid b)$, система \textbf{несовместна} (решений нет).
\end{enumerate}
\end{theorem}

\subsection{Почему элементарные преобразования строк не изменяют решений?}
\begin{definition}[Элементарные строчные операции]
Существуют три типа элементарных преобразований строк матрицы:
\begin{enumerate}
    \item Перестановка двух строк местами ($R_i \leftrightarrow R_j$);
    \item Умножение строки на ненулевой скаляр ($R_i \to c R_i$, $c \ne 0$);
    \item Прибавление к одной строке другой строки, умноженной на скаляр ($R_i \to R_i + c R_j$).
\end{enumerate}
\end{definition}

\begin{theorem}[О сохранении множества решений]
Элементарные преобразования строк расширенной матрицы $(A \mid b)$ не изменяют множества решений системы $Ax = b$.
\end{theorem}
\begin{proof}
    Любое элементарное преобразование строк эквивалентно умножению расширенной матрицы слева на соответствующую
    \textbf{элементарную матрицу} $E \in \mathbb{R}^{m \times m}$:
    \[ E(A \mid b) = (EA \mid Eb). \]
    Фундаментальное свойство элементарных матриц заключается в том, что \textbf{все они строго обратимы},
    причём матрица $E^{-1}$ сама является элементарной матрицей того же типа:
    \begin{itemize}
        \item Матрица перестановки строк $P_{ij}$ совпадает со своей обратной: $P_{ij}^{-1} = P_{ij}$;
        \item Обратная к умножению строки на $c \ne 0$ есть умножение на $\frac{1}{c}$;
        \item Обратная к прибавлению строки $j$, умноженной на $c$, есть вычитание строки $j$, умноженной на $c$.
    \end{itemize}
    Пусть $x$ — решение исходной системы, то есть $Ax = b$.
    Умножив обе части равенства слева на $E$, получаем $E(Ax) = Eb \iff (EA)x = Eb$.
    Следовательно, $x$ является решением преобразованной системы.
    
    Обратно, пусть $x$ — решение преобразованной системы, то есть $(EA)x = Eb$.
    Поскольку существует обратная матрица $E^{-1}$, умножим равенство слева на $E^{-1}$:
    \[ E^{-1}((EA)x) = E^{-1}(Eb) \iff (E^{-1}E)Ax = (E^{-1}E)b \iff Ax = b. \]
    Таким образом, $x$ является решением исходной системы.
    Множества решений совпадают.
\end{proof}

\section{LU-разложение и его вычислительное преимущество}

\begin{definition}[LU-decomposition]
\textbf{$LU$-разложением} квадратной матрицы $A \in \mathbb{R}^{n \times n}$ называется представление вида:
\[ A = L U, \]
где $L \in \mathbb{R}^{n \times n}$ — нижнетреугольная матрица (с единицами на главной диагонали в случае разложения
Дулиттла), а $U \in \mathbb{R}^{n \times n}$ — верхнетреугольная матрица.
Если в процессе исключения Гаусса требуются перестановки строк (выбор главного элемента для численной устойчивости),
используется разложение с перестановкой:
\[ P A = L U, \]
где $P$ — матрица перестановки.
\end{definition}

\subsection{Вычислительное преимущество LU-разложения}
Пусть требуется решить систему $Ax = b$.
Подставив $A = LU$, получаем:
\[ L (U x) = b. \]
Решение распадается на два последовательных простых треугольных этапа:
\begin{enumerate}
    \item \textbf{Прямая подстановка (Forward substitution):} решаем систему с нижнетреугольной матрицей $L y = b$.
    Требует $\approx \frac{n^2}{2}$ операций.
    \item \textbf{Обратная подстановка (Back substitution):} решаем систему с верхнетреугольной матрицей $U x = y$.
    Требует $\approx \frac{n^2}{2}$ операций.
\end{enumerate}

\textbf{Сравнение сложности:}
\begin{itemize}
    \item Прямой метод Гаусса «с нуля» требует $\mathcal{O}\left(\frac{2}{3} n^3\right)$ арифметических операций
    для каждого нового вектора правой части $b$.
    \item Факторизация $A=LU$ вычисляется \textbf{один раз} со сложностью $\mathcal{O}\left(\frac{2}{3} n^3\right)$.
    \item После этого для любого нового вектора $b$ нахождение решения через $Ly = b$ и $Ux = y$ требует всего лишь
    $\mathcal{O}(n^2)$ операций!
\end{itemize}
В реальных прикладных задачах (симуляция динамических систем, нелинейные итерации метода Ньютона, обработка
сигналов), где матрица $A$ неизменна, а правая часть $b$ меняется тысячи раз за секунду, использование
$LU$-разложения ускоряет расчёты в тысячи раз.

\section{Типичные ошибки размерностей при работе с матрицами}

\begin{enumerate}
    \item \textbf{Несогласованность внутренних размерностей при умножении:} Матричное произведение $A B$
    определено тогда и только тогда, когда число столбцов левой матрицы совпадает с числом строк правой матрицы:
    \[ A_{(m \times k)} \cdot B_{(k \times n)} = C_{(m \times n)}. \]
    Попытка умножить $A_{(m \times k)}$ на $B_{(p \times n)}$ при $k \ne p$ является грубейшей синтаксической и
    содержательной ошибкой.
    \item \textbf{Путаница между столбцами ($n \times 1$) и строками ($1 \times n$):} В математике по умолчанию
    вектор $x \in \mathbb{R}^n$ является столбцом ($n \times 1$).
    \begin{itemize}
        \item $x^\top y$ — \textbf{внутреннее (скалярное) произведение}: матрица
        $(1 \times n) \times (n \times 1) = (1 \times 1) \implies \text{число}$;
        \item $x y^\top$ — \textbf{внешнее произведение (outer product)}: матрица
        $(n \times 1) \times (1 \times m) = (n \times m) \implies \text{матрица ранга 1}$.
    \end{itemize}
    \item \textbf{Ошибки транспонирования произведения:}
    \[ (A B)^\top = B^\top A^\top \quad (\text{порядок сомножителей обязательно обращается!}). \]
    Написание $(AB)^\top = A^\top B^\top$ в общем случае приводит даже к несовпадению размерностей.
    \item \textbf{Ложное предположение о коммутативности:} В общем случае $A B \ne B A$.
    Более того, если $A B$ определено, произведение $B A$ может даже не существовать из-за несовпадения
    размерностей.
\end{enumerate}

\newpage

% =========================================================================
% ГЛАВА 3
% =========================================================================

    \chapter{Собственные значения, спектральная теория и матричные разложения}
        \section{Собственные векторы и характеристическое уравнение}

    \begin{definition}[Собственное значение и собственный вектор]
Пусть $A \in \mathbb{R}^{n \times n}$ — квадратная матрица.
Ненулевой вектор $v \in \mathbb{C}^n$ ($v \ne 0$) называется \textbf{собственным вектором} матрицы $A$, а число $\lambda
\in \mathbb{C}$ — соответствующим ему \textbf{собственным значением}, если выполнено соотношение:
\[ A v = \lambda v. \]
    \end{definition}

    \begin{theorem}[Вывод характеристического уравнения]
Равенство $Av = \lambda v$ эквивалентно однородной системе:
\[ (A - \lambda I_n) v = 0. \]
По определению собственного вектора требуется найти \textbf{ненулевое} решение $v \ne 0$.
Ненулевое решение однородной системы существует тогда и только тогда, когда ядро матрицы $(A - \lambda I_n)$
нетривиально:
\[ \ker(A - \lambda I_n) \ne \{0\} \iff \rank(A - \lambda I_n) < n. \]
В силу критерия вырожденности квадратной матрицы это эквивалентно равенству нулю ее определителя:
\[ p_A(\lambda) \coloneqq \det(A - \lambda I_n) = 0. \]
Уравнение $\det(A - \lambda I_n) = 0$ называется \textbf{характеристическим уравнением} матрицы $A$, а $p_A(\lambda)$ —
\textbf{характеристическим многочленом} степени $n$.
    \end{theorem}

    \section{Алгебраическая и геометрическая кратности}

    \begin{definition}[Кратности собственных значений]
Пусть $\lambda_0$ — собственное значение матрицы $A$.
\begin{enumerate}
    \item \textbf{Алгебраическая кратность} $\mu_A(\lambda_0)$ (или $\operatorname{alg\_mult}(\lambda_0)$) — это
    кратность корня $\lambda_0$ в характеристическом многочлене $p_A(\lambda)$:
    \[ p_A(\lambda) = (\lambda - \lambda_0)^{\mu_A(\lambda_0)} q(\lambda), \quad q(\lambda_0) \ne 0. \]
    \item \textbf{Геометрическая кратность} $\gamma_A(\lambda_0)$ (или $\operatorname{geom\_mult}(\lambda_0)$) — это
    размерность соответствующего собственного подпространства $E_{\lambda_0} = \ker(A - \lambda_0 I_n)$:
    \[ \gamma_A(\lambda_0) \coloneqq \dim \ker(A - \lambda_0 I_n) = n - \rank(A - \lambda_0 I_n). \]
\end{enumerate}
    \end{definition}

    \begin{theorem}[Неравенство кратностей]
Для любого собственного значения $\lambda_0$ геометрическая кратность всегда не превосходит алгебраическую:
\[ 1 \le \gamma_A(\lambda_0) \le \mu_A(\lambda_0). \]
Матрица $A \in \mathbb{R}^{n \times n}$ диагонализируема тогда и только тогда, когда все корни характеристического
многочлена лежат в поле и для каждого собственного значения геометрическая кратность совпадает с алгебраической:
$\gamma_A(\lambda_i) = \mu_A(\lambda_i)$.
    \end{theorem}

    \section{Линейная независимость собственных векторов для различных собственных значений}

    \begin{theorem}
        Собственные векторы $v_1, v_2, \dots, v_k$, отвечающие попарно различным собственным значениям $\lambda_1,
        \lambda_2, \dots, \lambda_k$ ($\lambda_i \ne \lambda_j$ при $i \ne j$), линейно независимы.
    \end{theorem}
    \begin{proof}
        Проведем строгое доказательство методом математической индукции по числу векторов $k$.

        \textbf{База индукции ($k = 1$):} По определению собственный вектор $v_1 \ne 0$.
        Любая система из одного ненулевого вектора линейно независима ($c_1 v_1 = 0 \implies c_1 = 0$).
        База доказана.

        \textbf{Индукционный переход:} Предположим, что теорема справедлива для любой системы из $k - 1$ собственных
        векторов, отвечающих различным собственным значениям, то есть система $\{v_1, \dots, v_{k-1}\}$ линейно
        независима.

        Докажем независимость системы $\{v_1, \dots, v_k\}$.
        Рассмотрим тривиальную линейную комбинацию:
        \begin{equation} \label{eq:ind1} c_1 v_1 + c_2 v_2 + \dots + c_{k-1} v_{k-1} + c_k v_k = 0.
        \end{equation}
        Применим оператор $A$ к обеим частям равенства \eqref{eq:ind1}.
        В силу линейности оператора $A$ и равенств $A v_i = \lambda_i v_i$:
        \begin{equation} \label{eq:ind2}
            c_1 \lambda_1 v_1 + c_2 \lambda_2 v_2 + \dots + c_{k-1} \lambda_{k-1} v_{k-1} + c_k \lambda_k v_k = 0.
        \end{equation}
        Теперь умножим исходное равенство \eqref{eq:ind1} на скаляр $\lambda_k$:
        \begin{equation} \label{eq:ind3}
            c_1 \lambda_k v_1 + c_2 \lambda_k v_2 + \dots + c_{k-1} \lambda_k v_{k-1} + c_k \lambda_k v_k = 0.
        \end{equation}
        Вычтем равенство \eqref{eq:ind3} из равенства \eqref{eq:ind2}:
        \[ c_1 (\lambda_1 - \lambda_k) v_1 + c_2 (\lambda_2 - \lambda_k) v_2 + \dots + c_{k-1} (\lambda_{k-1} -
        \lambda_k) v_{k-1} + 0 = 0. \]
        Мы получили линейную комбинацию первых $k - 1$ векторов.
        По индукционному предположению система $\{v_1, \dots, v_{k-1}\}$ линейно независима.
        Следовательно, все коэффициенты этой комбинации обязаны равняться нулю:
        \[ c_i (\lambda_i - \lambda_k) = 0 \quad \forall i \in \{1, 2, \dots, k-1\}. \]
        Так как по условию все собственные значения попарно различны, разность $(\lambda_i - \lambda_k) \ne 0$ для всех
        $i < k$.
        Отсюда немедленно следует:
        \[ c_i = 0 \quad \forall i \in \{1, 2, \dots, k-1\}. \]
        Подставляя $c_1 = c_2 = \dots = c_{k-1} = 0$ в исходное равенство \eqref{eq:ind1}, получаем $c_k v_k = 0$.
        А так как $v_k \ne 0$, получаем $c_k = 0$.
        Таким образом, все коэффициенты $c_1 = c_2 = \dots = c_k = 0$, что доказывает линейную независимость системы
        $\{v_1, \dots, v_k\}$.
    \end{proof}

    \section{Диагонализация и возведение матриц в степень}

    \begin{theorem}[Диагонализация через собственный базис]
Пусть матрица $A \in \mathbb{R}^{n \times n}$ имеет $n$ линейно независимых собственных векторов $v_1, \dots, v_n$,
отвечающих собственным значениям $\lambda_1, \dots, \lambda_n$.
Сформируем матрицу $P$, столбцами которой являются собственные векторы, и диагональную матрицу $\Lambda$:
\[
    P = \begin{pmatrix}
            | & & | \\ v_1 & \dots & v_n \\ | & & |
    \end{pmatrix}, \quad \Lambda = \begin{pmatrix}
                                       \lambda_1 & &
                                       0 \\ & \ddots & \\ 0 & & \lambda_n
    \end{pmatrix}.
\]
Тогда:
\[ A P = P \Lambda \iff A = P \Lambda P^{-1}. \]
    \end{theorem}

    \subsection{Вычисление степеней матриц}
    Если матрица диагонализируема, возведение ее в степень $k \in \mathbb{N}$ тривиализуется:
    \[ A^k=(P \Lambda P^{-1})^k = (P \Lambda P^{-1})(P \Lambda P^{-1}) \dots (P \Lambda P^{-1}) = P \Lambda^k P^{-1}. \]
    Возведение диагональной матрицы $\Lambda$ в степень сводится к возведению ее диагональных элементов в степень $k$:
    \[ \Lambda^k = \diag(\lambda_1^k, \lambda_2^k, \dots, \lambda_n^k). \]
    Сложность прямого умножения матриц составляет $\mathcal{O}(k n^3)$, в то время как через диагонализацию она
    составляет $\mathcal{O}(n^3 + n \log k)$!
    Это позволяет мгновенно исследовать асимптотическое поведение динамических систем, цепей Маркова и разностных
    уравнений при $k \to \infty$.

    \section{Вещественная спектральная теорема для симметричных матриц}

    \begin{theorem}[Вещественная спектральная теорема]
Пусть $A \in \mathbb{R}^{n \times n}$ — вещественная симметричная матрица ($A^\top = A$).
Тогда:
\begin{enumerate}
    \item Все собственные значения матрицы $A$ вещественны: $\lambda_i \in \mathbb{R}$;
    \item Собственные векторы, отвечающие различным собственным значениям, взаимно ортогональны;
    \item Существует \textbf{ортонормированный базис} пространства $\mathbb{R}^n$, состоящий из собственных векторов
    матрицы $A$.
\end{enumerate}
Следовательно, матрица $A$ ортогонально диагонализируема:
\[ A = Q \Lambda Q^\top = Q \Lambda Q^{-1}, \]
где $Q \in \mathbb{R}^{n \times n}$ — ортогональная матрица ($Q^\top Q = Q Q^\top = I$), а $\Lambda = \diag(\lambda_1,
\dots, \lambda_n)$.
    \end{theorem}

    \begin{remark}[Спектральное разложение симметричной матрицы]
Записав ортогональную диагонализацию в блочной форме через столбцы $Q = [q_1 \dots q_n]$, получаем \textbf{спектральное
разложение}:
\[ A = \sum_{i=1}^n \lambda_i q_i q_i^\top. \]
Каждое слагаемое $q_i q_i^\top$ является матрицей ортогонального проектора ранга 1 на одномерное собственное
подпространство $\spanop(q_i)$.
    \end{remark}

    \section{Положительно определенные матрицы, Cholesky и QR}

    \begin{definition}[Положительно определенная матрица]
Симметричная матрица $A \in \mathbb{R}^{n \times n}$ называется \textbf{положительно определённой} (обозначается $A
\succ 0$), если соответствующая квадратичная форма строго положительна для любого ненулевого вектора:
\[ x^\top A x > 0 \quad \forall x \in \mathbb{R}^n \setminus \{0\}. \]
    \end{definition}

    \begin{theorem}[Критерии положительной определённости]
Для симметричной матрицы $A = A^\top$ следующие утверждения эквивалентны:
\begin{enumerate}
    \item $A \succ 0$;
    \item Все собственные значения строго положительны: $\lambda_i(A) > 0$ для всех $i = 1, \dots, n$;
    \item \textbf{Критерий Сильвестра:} Все угловые главные миноры строго положительны ($\Delta_k > 0$ для
    $k = 1, \dots, n$);
    \item Существует единственная невырожденная нижнетреугольная матрица $L$ с положительными элементами на
    диагонали такая, что: \[ A = L L^\top \quad \text{(\textbf{Разложение Холецкого})}. \]
\end{enumerate}
    \end{theorem}

    \begin{definition}[QR-decomposition]
\textbf{$QR$-разложением} матрицы $A \in \mathbb{R}^{m \times n}$ ($m \ge n$) называется факторизация:
\[ A = Q R, \]
где $Q \in \mathbb{R}^{m \times n}$ имеет ортонормированные столбцы ($Q^\top Q = I_n$), а $R \in \mathbb{R}^{n \times
n}$ — верхнетреугольная невырожденная матрица с положительными диагональными элементами.
Разложение получается в результате ортогонализации Грама-Шмидта столбцов матрицы $A$.
    \end{definition}

    \section{Сравнительный анализ матричных разложений}

    В таблице~\ref{tab:decompositions} приведен детальный сравнительный обзор основных разложений курса.

    \begin{table}[htbp]
        \centering
        \small
        \begin{tabular}{lp{3.5cm}p{3cm}cp{4cm}}
            \toprule
            \textbf{Разложение} & \textbf{Требования к матрице} & \textbf{Вид сомножителей} & \textbf{Сложность} &
            \textbf{Основное применение} \\ \midrule
            $LU$ & Квадратная $n \times n$, обратима (с перестановками $P$) & $L$ нижнетреугольная, $U$
            верхнетреугольная & $\frac{2}{3} n^3$ & Решение СЛУ $Ax=b$, вычисление определителя и обратной матрицы \\
            \midrule
            Холецкий & Квадратная $n \times n$, $A = A^\top \succ 0$ & $A = L L^\top$, $L$ нижнетреугольная &
            $\frac{1}{3} n^3$ & Оптимизация, МНК, многомерные нормальные распределения (в 2 раза быстрее LU) \\
            \midrule
            $QR$ & Произвольная $m \times n$, $m \ge n$ & $Q$ ортонормированная, $R$ верхнетреугольная
            & $2 n^2 \left(m - \frac{n}{3}\right)$ & Задача МНК, численный алгоритм поиска собственных значений
            (QR-алгоритм) \\ \midrule
            Спектральное & Квадратная $n \times n$, $A = A^\top$
            & $A = Q \Lambda Q^\top$, $Q$ ортогональная, $\Lambda$ диагональная & $\mathcal{O}(n^3)$ & Анализ
            квадратичных форм, PCA, физические колебания, динамические системы \\ \midrule
            $SVD$ & \textbf{Любая} прямоугольная $m \times n$ & $U, V$ ортогональные, $\Sigma$ псевдодиагональная
            & $\mathcal{O}(m n^2)$
            & Низкоранговая аппроксимация, сжатие данных, псевдообращение, анализ главных направлений \\ \bottomrule
        \end{tabular}
        \caption{Сводная классификация фундаментальных матричных разложений}
        \label{tab:decompositions}
    \end{table}

    \newpage

% =========================================================================
% ГЛАВА 4
% =========================================================================

    \chapter{Сингулярное разложение (SVD) и прикладные методы}
        \section{Теорема о сингулярном разложении}

    \begin{theorem}[О сингулярном разложении (SVD)]
Для любой вещественной матрицы $A \in \mathbb{R}^{m \times n}$ ранга $r \le \min(m, n)$ существуют ортогональные матрицы
$U \in \mathbb{R}^{m \times m}$ ($U^\top U = I_m$) и $V \in \mathbb{R}^{n \times n}$ ($V^\top V = I_n$), а также
диагонально-подобная матрица $\Sigma \in \mathbb{R}^{m \times n}$ такие, что:
\[ A = U \Sigma V^\top. \]
При этом матрица $\Sigma$ имеет вид:
\[
    \Sigma = \begin{pmatrix}
                 \sigma_1 & & & 0 & \dots & 0 \\
                 & \ddots & & \vdots & \ddots & \vdots \\
                 & & \sigma_r & 0 & \dots & 0 \\
                 0 & \dots & 0 & 0 & \dots & 0 \\
                 \vdots & \ddots & \vdots & \vdots & \ddots & \vdots \\
                 0 & \dots & 0 & 0 & \dots & 0
    \end{pmatrix},
\]
где числа $\sigma_1 \ge \sigma_2 \ge \dots \ge \sigma_r > 0$ называются \textbf{сингулярными числами} матрицы $A$,
столбцы $u_1, \dots, u_m$ матрицы $U$ — \textbf{левыми сингулярными векторами}, а столбцы $v_1, \dots, v_n$ матрицы $V$
— \textbf{правыми сингулярными векторами}.
    \end{theorem}

    \section{Конструктивный вывод SVD из произведения $A^\top A$}

    Рассмотрим матрицу $S \coloneqq A^\top A \in \mathbb{R}^{n \times n}$.
    \begin{enumerate}
        \item \textbf{Симметричность:} $S^\top = (A^\top A)^\top = A^\top (A^\top)^\top = A^\top A = S$.
        \item \textbf{Неотрицательная определенность:} Для любого $x \in \mathbb{R}^n$:
        \[ x^\top S x = x^\top (A^\top A) x = (Ax)^\top (Ax) = \|Ax\|_2^2 \ge 0. \]
        Следовательно, по спектральной теореме все собственные значения матрицы $A^\top A$ вещественны и неотрицательны:
        \[ \lambda_1 \ge \lambda_2 \ge \dots \ge \lambda_r > \lambda_{r+1} = \dots = \lambda_n = 0, \quad \text{где }
        r = \rank(A^\top A) = \rank(A). \]
        Существует ортонормированный базис $\mathbb{R}^n$ из собственных векторов $(v_1,\dots, v_n)$ матрицы $A^\top A$:
        \[ A^\top A v_i = \lambda_i v_i, \quad v_i^\top v_j = \delta_{ij}. \]
        Векторы $v_1, \dots, v_n$ объявляются \textbf{правыми сингулярными векторами}, а \textbf{сингулярные числа}
        определяются формулой:
        \[ \sigma_i \coloneqq \sqrt{\lambda_i} \ge 0. \]
        \item \textbf{Построение левых сингулярных векторов:} Для $i \in \{1, 2, \dots, r\}$ определим векторы:
        \[ u_i \coloneqq \frac{1}{\sigma_i} A v_i \in \mathbb{R}^m. \]
        Проверим их взаимную ортонормированность:
        \[ \langle u_i, u_j \rangle = u_i^\top u_j = \left(\frac{1}{\sigma_i} A v_i\right)^\top
        \left(\frac{1}{\sigma_j} A v_j\right) = \frac{1}{\sigma_i \sigma_j} v_i^\top (A^\top A v_j) =
        \frac{1}{\sigma_i \sigma_j} v_i^\top (\lambda_j v_j) = \frac{\lambda_j}{\sigma_i \sigma_j} (v_i^\top v_j). \]
        Учитывая, что $v_i^\top v_j = \delta_{ij}$ и $\lambda_i = \sigma_i^2$:
        \[ \langle u_i, u_j \rangle = \frac{\sigma_j^2}{\sigma_i \sigma_j} \delta_{ij} = \delta_{ij}. \]
        Таким образом, система векторов $(u_1, \dots, u_r)$ строго ортонормирована в $\mathbb{R}^m$.
        \item Дополним систему $(u_1, \dots, u_r)$ по Граму-Шмидту до полного ортонормированного базиса
        $(u_1, \dots, u_m)$ пространства $\mathbb{R}^m$.
        \item Заметим, что по определению $A v_i = \sigma_i u_i$ при $i \le r$, и $A v_i = 0$ при $i > r$
        (так как $\|A v_i\|^2 = \lambda_i = 0$).
        В матричной форме это записывается как: \[ A V = U \Sigma \implies A = U \Sigma V^\top. \]
    \end{enumerate}

    \section{Теорема об общих ненулевых собственных значениях $A^\top A$ и $A A^\top$}

    \begin{theorem}
        Матрицы $A^\top A \in \mathbb{R}^{n \times n}$ и $A A^\top \in \mathbb{R}^{m \times m}$ имеют в точности
        одни и те же ненулевые собственные значения с одинаковыми кратностями.
    \end{theorem}
    \begin{proof}
        Пусть $\lambda \ne 0$ — ненулевое собственное значение матрицы $A^\top A$, и
        $v \in \mathbb{R}^n$ — соответствующий собственный вектор:
        \[ (A^\top A) v = \lambda v, \quad v \ne 0. \]
        Умножим обе части равенства слева на матрицу $A$:
        \[ A (A^\top A v) = A (\lambda v) \iff (A A^\top) (A v) = \lambda (A v). \]
        Обозначим вектор $w \coloneqq A v \in \mathbb{R}^m$.
        Убедимся, что $w \ne 0$.
        Вычислим его евклидову норму в квадрате:
        \[ \|w\|_2^2 = w^\top w = (Av)^\top (Av) = v^\top A^\top A v = v^\top (\lambda v) = \lambda \|v\|_2^2. \]
        Так как $\lambda > 0$ и $v \ne 0$, получаем $\|w\|_2^2 > 0 \implies w \ne 0$.
        Следовательно, вектор $w = Av$ является ненулевым собственным вектором матрицы $A A^\top$, отвечающим тому
        же самому собственному значению $\lambda$.

        В силу симметрии (заменяя $A$ на $A^\top$), аналогично доказывается, что любое ненулевое собственное
        значение матрицы $A A^\top$ является собственным значением матрицы $A^\top A$.
    \end{proof}

    \section{Формы SVD: полное, редуцированное и внешнее произведение}

    \begin{enumerate}
        \item \textbf{Полное SVD (Full SVD):}
        \[ A = U \Sigma V^\top, \quad U \in \mathbb{R}^{m \times m}, \; \Sigma \in \mathbb{R}^{m \times n}, \; V \in
        \mathbb{R}^{n \times n}. \]
        \item \textbf{Редуцированное / компактное SVD (Reduced / Thin SVD):}
        Исключаются нулевые блоки сингулярных чисел.
        Сохраняются только $r = \rank(A)$ ненулевых компонент:
        \[ A = U_r \Sigma_r V_r^\top, \quad U_r \in \mathbb{R}^{m \times r}, \; \Sigma_r =
        \diag(\sigma_1, \dots, \sigma_r) \in \mathbb{R}^{r \times r}, \; V_r \in \mathbb{R}^{n \times r}. \]
        Столбцы $U_r$ и $V_r$ ортонормированы: $U_r^\top U_r = I_r$, $V_r^\top V_r = I_r$.
        \item \textbf{Внешне-произведённая форма (Outer-Product Form):}
        Используя блочное правило умножения «столбец на строку»:
        \[ A = \sum_{i=1}^r \sigma_i u_i v_i^\top. \]
        Каждое слагаемое $u_i v_i^\top \in \mathbb{R}^{m \times n}$ представляет собой элементарный строительный
        кирпичик — матрицу ранга 1, а сингулярное число $\sigma_i$ выступает в роли «веса» или «силы сигнала»
        данного канала передачи информации.
    \end{enumerate}

    \section{Геометрический смысл SVD}

    Линейное отображение, задаваемое любой вещественной матрицей $A \in \mathbb{R}^{m \times n}$, переводит единичную
    гиперсферу пространства $\mathbb{R}^n$ в гиперэллипсоид пространства $\mathbb{R}^m$.
    Действие $A x = U \Sigma V^\top x$ концептуально распадается на 3 последовательных геометрических шага:
    \begin{enumerate}
        \item \textbf{Ортогональный поворот/отражение в пространстве прообразов ($V^\top$):} Вектор $x$
        поворачивается в базис правых сингулярных векторов $v_i$.
        Сфера переходит в сферу.
        \item \textbf{Масштабирование вдоль координатных осей ($\Sigma$):} Пространство растягивается вдоль $i$-й оси с
        коэффициентом $\sigma_i$.
        Сфера деформируется в эллипсоид, полуоси которого равны $\sigma_i$.
        \item \textbf{Ортогональный поворот/отражение в пространстве образов ($U$):} Полученный эллипсоид поворачивается
        в пространстве $\mathbb{R}^m$ так, что направления его главных полуосей совпадают с левыми сингулярными
        векторами $u_i$.
    \end{enumerate}

    \section{Низкоранговая аппроксимация и теорема Эккарта-Янга-Мирского}

    \begin{definition}[Усечённое SVD]
Для ранга аппроксимации $k < r = \rank(A)$ \textbf{усечённым сингулярным разложением} называется матрица:
\[ A_k \coloneqq \sum_{i=1}^k \sigma_i u_i v_i^\top = U_k \Sigma_k V_k^\top. \]
Очевидно, что $\rank(A_k) = k$.
    \end{definition}

    \begin{theorem}[Теорема Эккарта-Янга-Мирского]
Матрица $A_k$ является наилучшей аппроксимацией матрицы $A$ среди всех матриц ранга не выше $k$ как в спектральной
норме, так и в норме Фробениуса:
\begin{gather*}
    \min_{\rank(B) \le k} \|A - B\|_2 = \|A - A_k\|_2 = \sigma_{k+1},\\
    \min_{\rank(B) \le k} \|A - B\|_F = \|A - A_k\|_F = \sqrt{\sum_{i=k+1}^r \sigma_i^2}.\\
\end{gather*}
    \end{theorem}

    \begin{example}[Применения низкоранговой аппроксимации]
\begin{itemize}
    \item \textbf{Сжатие изображений:} Матрица пикселей размера $m \times n$ требует хранения $m \cdot n$ чисел.
    Усечённое разложение $A_k$ хранит лишь $k(m + n + 1)$ чисел.
    При $k \ll \min(m, n)$ достигается многократное сжатие с сохранением основных визуальных деталей.
    \item \textbf{Латентно-семантический анализ (LSA):} Матрица «термы-документы» проецируется на $k$ скрытых
    семантических тем, фильтруя синонимический шум.
\end{itemize}
    \end{example}

    \section{Псевдообратная матрица Мура-Пенроуза}

    \begin{definition}[Псевдообратная матрица Мура-Пенроуза]
Пусть $A = U \Sigma V^\top$ — сингулярное разложение матрицы $A \in \mathbb{R}^{m \times n}$.
\textbf{Псевдообратной матрицей Мура-Пенроуза} называется матрица $A^+ \in \mathbb{R}^{n \times m}$, определяемая
формулой:
\[ A^+ \coloneqq V \Sigma^+ U^\top = \sum_{i=1}^r \frac{1}{\sigma_i} v_i u_i^\top, \]
где матрица $\Sigma^+ \in \mathbb{R}^{n \times m}$ получается из $\Sigma^\top$ заменой каждого ненулевого сингулярного
числа $\sigma_i > 0$ на обратную величину $\frac{1}{\sigma_i}$, а нули остаются на местах.
    \end{definition}

    \begin{theorem}[Условия Мура-Пенроуза]
Матрица $A^+$ является единственной матрицей, удовлетворяющей 4 условиям Мура-Пенроуза:
\begin{enumerate}
    \item $A A^+ A = A$;
    \item $A^+ A A^+ = A^+$;
    \item $(A A^+)^\top = A A^+$ (ортогональный проектор на $\Col(A)$);
    \item $(A^+ A)^\top = A^+ A$ (ортогональный проектор на $\Row(A)$).
\end{enumerate}
    \end{theorem}

    \begin{theorem}[Оптимальность псевдообратного решения]
Для любой системы $Ax = b$ вектор $x^* \coloneqq A^+ b$:
\begin{enumerate}
    \item Минимизирует невязку наименьших квадратов $\|Ax - b\|_2$:
    \[ \|A x^* - b\|_2 = \min_{x \in \mathbb{R}^n} \|A x - b\|_2. \]
    \item Если решений задачи наименьших квадратов бесконечно много, $x^*$ является решением с \textbf{минимальной
    евклидовой нормой}:
    \[ \|x^*\|_2 = \min \{ \|x\|_2 \mid x \in \operatorname{Argmin} \|Ax - b\|_2 \}. \]
\end{enumerate}
    \end{theorem}

    \section{Рекомендательные системы и латентный факторный анализ}

    Пусть $R \in \mathbb{R}^{m \times n}$ — частично заполненная матрица рейтингов, где строка $u$ соответствует
    пользователю, а столбец $i$ — фильму или товару.
    Используя усечённое SVD ранга $k$:
    \[ R \approx U_k \Sigma_k V_k^\top = (U_k \Sigma_k^{1/2}) (V_k \Sigma_k^{1/2})^\top = P Q^\top. \]
    \begin{itemize}
        \item Строки матрицы $P = U_k \Sigma_k^{1/2} \in \mathbb{R}^{m \times k}$ являются $k$-мерными векторами скрытых
        предпочтений пользователей $p_u$;
        \item Строки матрицы $Q = V_k \Sigma_k^{1/2} \in \mathbb{R}^{n \times k}$ являются $k$-мерными векторами 
        латентных характеристик элементов $q_i$;
        \item Прогноз рейтинга для не просмотренного пользователем фильма вычисляется через скалярное произведение 
        латентных  векторов: \[ \hat{r}_{ui} = \langle p_u, q_i \rangle = \sum_{f=1}^k p_{uf} q_{if}. \]
    \end{itemize}

\end{document}
